\section{Fase 4: protocolo de evaluación}
\label{sec:fase4}

Se emplean dos protocolos de validación, que responden a preguntas distintas. El
protocolo intra-sujeto mide cuánta información motora puede decodificarse en un
sujeto a partir de sus propios datos de entrenamiento. El protocolo inter-sujeto
mide cuánto de ese aprendizaje se transfiere a un sujeto nuevo. La comparación
entre ambos es la que motiva el indicador de la sección~\ref{sec:dsi}.

\subsection{Validación intra-sujeto}
\label{sec:within}

Cada sujeto se evalúa consigo mismo y \textbf{por separado en cada sesión}, para
que el resultado no mezcle la variabilidad entre días con la del propio
protocolo. Los 200 ensayos de una sesión se reparten en diez \textit{folds} estratificados por
clase, sorteados con barajado dentro de cada sesión (\texttt{StratifiedKFold}), y no en
bloques temporales contiguos; en cada partición un \textit{fold} distinto actúa como prueba y los
nueve restantes se usan para entrenar y validar, como muestra la
figura~\ref{fig:within}.

\begin{figure}[H]
\centering
\begin{tikzpicture}[
  font=\footnotesize,
  bloque/.style={draw, minimum width=0.82cm, minimum height=0.75cm},
  fl/.style={-{Latex[length=1.8mm]}, semithick},
  x=1cm, y=1cm
]
\foreach \i in {1,...,10}{
  \node[bloque, fill=blue!8] at ({0.9*\i},0) {\scriptsize \i};
}
\node[bloque, fill=orange!25] (test) at ({0.9*3},0) {\scriptsize 3};

\node[anchor=south] (etq) at ({0.9*3},0.75) {\scriptsize prueba: 20 ensayos};
\draw[fl] (etq.south) -- (test.north);

\node[anchor=north, align=center] at (4.95,-0.55)
  {\scriptsize los otros nueve \textit{folds} (180 ensayos):
   144 de entrenamiento y 36 de validación, que determinan la parada temprana};
\node[anchor=north, align=center] at (4.95,-1.35)
  {\scriptsize el \textit{fold} de prueba rota $\rightarrow$ 10 particiones por sesión,
   2 sesiones por sujeto};
\end{tikzpicture}
\caption{Validación intra-sujeto. Los 200 ensayos de una sesión se dividen en
diez \textit{folds} estratificados por clase, sorteados al azar dentro de la sesión (no son
bloques temporales contiguos), y cada uno actúa como conjunto de prueba una vez. El proceso se
repite en las dos sesiones, lo que da 20 particiones por sujeto, y con tres semillas ($2026$, $2027$ y $2028$) en las
cuatro arquitecturas (sección~\ref{sec:semillas-por-arq}).}
\label{fig:within}
\end{figure}

Debe señalarse que en este protocolo cada modelo se entrena con
\textbf{144 ensayos} y se detiene con 36. Es un régimen de datos muy escaso que, junto
con el reducido tamaño de la validación (sección~\ref{sec:seleccion-intra}), condiciona
buena parte de lo que se observa en el capítulo~\ref{cap:resultados}.

\subsection{Validación \loso{}}

Se emplea validación \textit{Leave-One-Subject-Out} sobre los 54 sujetos. En cada
\textit{fold}, un sujeto constituye el conjunto de prueba y los $53$ restantes se reparten
en dos grupos: \textbf{dos sujetos completos} forman el conjunto de validación, cuya pérdida
determina la parada temprana, y los otros \textbf{51} forman el conjunto de entrenamiento. Los
dos sujetos de validación se eligen de forma determinista y rotan entre \textit{folds}, de
modo que ninguno actúa siempre como validación. El sujeto de prueba \textbf{no aporta ningún
ensayo} al entrenamiento, a la validación ni a los estadísticos de normalización.

Cada modelo se entrena así con $51 \times 360 = 18\,360$ ensayos, frente a los 144 del
protocolo anterior, y se detiene con $2 \times 400 = 800$ ensayos de personas que no
participan en el ajuste de los pesos. Que la parada temprana se decida sobre personas no
vistas tiene una consecuencia que conviene anotar: el modelo conservado es el que mejor
generaliza a personas nuevas dentro del conjunto de origen, y no el que mejor se ajusta a
las personas con las que se entrenó.

\subsection{Protocolos adicionales: modelo global y modelo con sujetos apartados}
\label{sec:protocolos-extra}

Los dos protocolos anteriores responden a preguntas dentro de \openbmi{}. Para evaluar el
cambio de dominio entre conjuntos de datos distintos hacen falta dos protocolos más, que se
describen aquí porque condicionan la lectura de todos los resultados del
capítulo~\ref{cap:validacion-externa}.

\subsubsection{Qué sesiones se emplean}

Conviene fijarlo antes que nada, porque no es simétrico entre los dos conjuntos y de ello
depende la interpretación de las cifras:

\begin{itemize}
  \item De \openbmi{} se emplean \textbf{sus dos sesiones}, en todos los protocolos. Cada
  sujeto aporta $400$ segmentos, $200$ por sesión, y los $54$ sujetos suman $21\,600$.
  \item De WBCIC-SHU se emplea \textbf{únicamente la primera de sus tres sesiones}, tanto
  para entrenar el modelo nativo como para evaluar la transferencia. Cada sujeto aporta
  $200$ segmentos. La versión distribuida contiene $10\,199$, uno menos que los $10\,200$
  nominales, y quedan $10\,198$ tras descartar un segmento saturado
  (sección~\ref{sec:wbcic-preproc}).
\end{itemize}

Restringirse a la primera sesión de WBCIC-SHU es una elección conservadora: sus autores
documentan un efecto de aprendizaje entre días: con \eegnet{}, la exactitud media sube de
$0.818$ en la primera sesión a $0.866$ en la segunda y $0.889$ en la tercera
\cite{yang2025wbcic}, de modo que la primera es la
condición menos favorable. La consecuencia, que se declara aquí y se retoma al interpretar
los resultados, es que el modelo nativo de WBCIC entrena con $144$ segmentos por persona
frente a los $360$ de un modelo de \openbmi{}: una sesión menos y, además, la fracción que
cede al conjunto reservado.

\subsubsection{Protocolo global: un solo modelo sobre todos los sujetos}

En los protocolos intra-sujeto y \loso{} el resultado es un promedio sobre decenas de
modelos: $20$ particiones por sujeto y semilla en el primero y $54$ modelos distintos en el segundo,
cada uno entrenado con $51$ sujetos y detenido con otros $2$. Esto impide trasladar \emph{un} modelo a otro conjunto de
datos, porque no existe un único conjunto de pesos que represente al protocolo.

El protocolo global entrena, por tanto, \textbf{un solo modelo por arquitectura} con la
totalidad de los sujetos disponibles, reservando un $10$\,\% estratificado por sujeto, clase
y sesión para la parada temprana. No produce ninguna medida de generalización por sí mismo:
su única salida son los pesos que después se aplican al segundo conjunto de datos. La
figura~\ref{fig:global} lo esquematiza.

\begin{figure}[H]
\centering
\begin{tikzpicture}[
  font=\footnotesize,
  caja/.style={draw, rounded corners=2pt, minimum height=0.8cm, align=center},
  fl/.style={-{Latex[length=1.8mm]}, semithick},
  x=1cm, y=1cm
]
\node[caja, fill=blue!10, minimum width=7.2cm] (tr) at (0,0)
  {\textbf{Entrenamiento}: 54 sujetos, 2 sesiones\\45 canales, $19\,440$ segmentos ($90$\,\%)};
\node[caja, fill=blue!22, minimum width=2.4cm] (va) at (4.9,0)
  {\textbf{Validación}\\$2\,160$ ($10$\,\%)\\[-2pt]\scriptsize parada temprana};
\node[anchor=north, font=\scriptsize, align=center] at (0,-0.85)
  {los mismos 54 sujetos en las dos cajas: la validación mide\\
   \textbf{sujetos vistos}, no generalización};
\node[caja, fill=orange!20, minimum width=4.4cm] (ck) at (0,-2.7)
  {Un único conjunto de pesos por arquitectura};
\draw[fl] (tr.south) ++(0,-1.2) -- (ck.north);
\node[caja, fill=green!14, minimum width=5.6cm] (ap) at (0,-4.2)
  {Aplicación a WBCIC-SHU \textbf{sin reentrenar}};
\draw[fl] (ck) -- (ap);
\end{tikzpicture}
\caption{Protocolo global. Se entrena un solo modelo por arquitectura con los 54 sujetos
para disponer de un conjunto de pesos transferible, con los 45 canales comunes a los dos
conjuntos. La validación determina la parada temprana, y su exactitud \textbf{no es una medida
de generalización}: procede de segmentos de los mismos sujetos que formaron parte del
entrenamiento.}
\label{fig:global}
\end{figure}

Este protocolo tiene una limitación que conviene enunciar de forma explícita, porque de
ella se deriva el siguiente: al entrenar con los $54$ sujetos \textbf{no queda ninguna persona
sin ver}. El modelo resultante no puede producir una cifra de \emph{persona nueva dentro
del mismo protocolo de adquisición}, y por tanto la caída que se observe al aplicarlo a otro
conjunto de datos mezcla dos cambios simultáneos, el de la persona y el del protocolo, sin
permitir separarlos.

Tampoco es lícito estimar el primero de esos dos efectos recurriendo al valor \loso{}. La
exactitud \loso{} es la media de $54$ modelos entrenados con $51$ sujetos cada uno, mientras
que la del protocolo global procede de un único modelo entrenado con $54$: son
\textbf{modelos distintos}, y su diferencia no aísla ningún factor.

\subsubsection{Protocolo con sujetos apartados: 44 más 10}

Para medir por separado el efecto de la persona y el del protocolo se entrena una segunda
tanda de modelos apartando \textbf{diez sujetos completos} de \openbmi{}, que no aportan
ningún segmento ni al entrenamiento ni a los estadísticos de normalización. Un mismo
conjunto de pesos se evalúa entonces sobre tres conjuntos, como muestra la
figura~\ref{fig:hold}. Del primero al segundo cambia la persona. Del segundo al tercero
cambia un \textbf{factor compuesto}: a la vez la población, la institución, el amplificador,
el montaje y las diferencias de preprocesamiento de la tabla~\ref{tab:armonizacion}, que
este diseño no permite separar.

\begin{figure}[H]
\centering
\begin{tikzpicture}[
  font=\footnotesize,
  caja/.style={draw, rounded corners=2pt, minimum height=0.8cm, align=center},
  ev/.style={draw, rounded corners=2pt, fill=green!14, minimum width=3.5cm,
             minimum height=1.35cm, align=center, font=\scriptsize},
  fl/.style={-{Latex[length=1.8mm]}, semithick},
  x=1cm, y=1cm
]
\node[caja, fill=blue!10, minimum width=5.6cm] (tr) at (-1.3,0)
  {\textbf{44 sujetos} de \openbmi{}\\45 canales, $15\,840$ segmentos};
\node[caja, fill=blue!22, minimum width=2.0cm] (va) at (2.6,0)
  {\textbf{Val.}\\$1\,760$};
\node[caja, fill=red!12, minimum width=2.6cm] (ho) at (5.6,0)
  {\textbf{10 sujetos}\\apartados};
\node[anchor=north, font=\scriptsize] at (5.6,-0.62) {nunca vistos};

\node[caja, fill=orange!20, minimum width=4.0cm] (ck) at (0,-1.9)
  {Un conjunto de pesos};
\draw[fl] (tr.south) -- (ck.north);

\node[ev] (e1) at (-4.6,-4.0)
  {\textbf{1. Validación}\\sujetos vistos\\determina la parada temprana};
\node[ev] (e2) at (0,-4.0)
  {\textbf{2. Los 10 apartados}\\\textbf{persona nueva}\\protocolo conocido};
\node[ev] (e3) at (4.6,-4.0)
  {\textbf{3. WBCIC-SHU}\\persona nueva\\\textbf{protocolo nuevo}};
\coordinate (bus) at (0,-2.9);
\draw[fl] (ck) -- (bus);
\foreach \e in {e1,e2,e3} \draw[fl] (bus) -| (\e.north);

\draw[{Latex[length=1.6mm]}-{Latex[length=1.6mm]}, dashed] (e1.south) ++(0,-0.25)
  -- node[below, font=\scriptsize] {efecto de la persona} ($(e2.south)+(0,-0.25)$);
\draw[{Latex[length=1.6mm]}-{Latex[length=1.6mm]}, dashed] (e2.south) ++(0,-0.95)
  -- node[below, font=\scriptsize] {efecto del protocolo} ($(e3.south)+(0,-0.95)$);
\end{tikzpicture}
\caption{Protocolo con sujetos apartados. Los tres conjuntos de evaluación se recorren con
\textbf{los mismos pesos}, de modo que las diferencias entre ellos no se deben a modelos
distintos: de 1 a 2 cambia la persona; de 2 a 3 cambia además el protocolo de adquisición,
que agrupa varios factores (tabla~\ref{tab:armonizacion}).
El conjunto 1 es también el que determina la parada temprana, y sus ensayos proceden de las
mismas sesiones que los de entrenamiento; ambas condiciones tienden a elevar su exactitud
y, con ella, la caída de 1 a 2.}
\label{fig:hold}
\end{figure}

La selección de los diez sujetos es aleatoria con semilla fija y se recalcula de forma
determinista a partir de la lista ordenada de sujetos, sin depender del orden en que el
sistema de ficheros los devuelva. Los sujetos apartados resultan ser
sub-04, 07, 13, 17, 30, 34, 36, 47, 51 y 52. Su exactitud \loso{} media, con la media de tres semillas por sujeto, queda por debajo
de la del resto en las cuatro arquitecturas, entre $2.2$ y $3.1$ puntos: con \eegnet{},
$0.742$ frente a $0.765$. La diferencia no es significativa (Mann-Whitney, $p = 0.56$ con
\eegnet{} y entre $0.34$ y $0.58$ con las otras tres), pero con diez sujetos la prueba no
tiene potencia para descartar un déficit de ese tamaño, que coincide con el margen de
equivalencia de la sección~\ref{sec:equivalencia}. Como el efecto de la persona se calcula
restando la exactitud sobre los apartados a la de los sujetos vistos, esa selección puede
sobrestimarlo en una magnitud comparable. Se hace
constar que \textbf{la semilla no se volvió a sortear tras comprobar esto}: repetir el
sorteo hasta obtener un subconjunto favorable sería una selección sesgada.

\subsubsection{Normalización del conjunto de destino}

Al aplicar un modelo a un conjunto de datos que no formó parte de su entrenamiento, la media y la desviación con que se
normaliza la entrada no están definidas de forma única. Se reportan tres alternativas, porque la
diferencia entre ellas es en sí misma informativa sobre la naturaleza del cambio de dominio:

\begin{itemize}
  \item \textbf{Origen}: se aplican la media y la desviación del entrenamiento en
  \openbmi{}. Es la condición más estricta, en la que no se emplea ninguna información del
  conjunto de destino.
  \item \textbf{Por segmento}: cada segmento se escala con una única media y una única
  desviación, calculadas sobre todos sus canales. No emplea ningún otro segmento, de modo
  que no es transductiva, y corrige una diferencia de ganancia común a todos los canales;
  no corrige diferencias de ganancia entre canales ni de covarianza espacial.
  \item \textbf{Por sujeto}: estadísticos de cada sujeto del destino. Sí es transductiva,
  pues supone disponer de todos los datos del sujeto antes de clasificar.
\end{itemize}

Esta elección permite contrastar una hipótesis concreta y acotada: \textbf{si la caída al
cambiar de conjunto se debiera a una diferencia de ganancia global entre los equipos de
registro, la normalización por segmento la corregiría y el desempeño aumentaría}. El
contraste tiene dos límites. No alcanza a diferencias de ganancia por canal ni de
covarianza. Y las tres políticas se aplican a un modelo entrenado con los estadísticos por
canal del origen, de modo que las dos alternativas introducen además un desajuste entre
la normalización de entrenamiento y la de prueba; un contraste completo exigiría entrenar
también con cada normalización. El capítulo~\ref{cap:validacion-externa} muestra que la
normalización por segmento no recupera el desempeño.

Finalmente, para que el modelo transferido y un modelo entrenado en el destino sean
comparables, se aparta en WBCIC-SHU un $20$\,\% de los segmentos de \textbf{cada} sujeto,
estratificado por sujeto y clase, con una semilla propia e independiente de la del
entrenamiento. Ese subconjunto, $2\,040$ segmentos, exactamente $40$ de cada uno de los
$51$ sujetos y equilibrado entre clases, no se emplea para entrenar ni para la parada
temprana de ninguno de los dos modelos, y es donde se miden ambos.

\subsection{Calibración con pocos ensayos del sujeto no visto}
\label{sec:metodo-fewshot}

El protocolo \loso{} mide el desempeño con cero ensayos del sujeto de prueba, y el
intra-sujeto con $144$ de entrenamiento. Para el escenario intermedio se calibra el modelo
con $k$ ensayos del sujeto excluido, con $k \in \{10, 20\}$. La calibración se ejecuta
dentro de cada \textit{fold} \loso{}, justo después de evaluar el modelo ya entrenado, sin
reentrenarlo, de modo que cada \textit{fold} y semilla tiene un único modelo de partida.

Los $400$ ensayos del sujeto excluido, de sus dos sesiones juntas, se dividen así. Para cada
$k$ y cada uno de $5$ sorteos se eligen $k$ ensayos de calibración, estratificados por clase
($k/2$ por clase), con un generador inicializado con $\mathrm{semilla}\times 1000 + k\times 10 +
\mathrm{sorteo}$. La evaluación es el resto del sujeto ($390$ ensayos si $k = 10$, $380$ si
$k = 20$); ningún ensayo de calibración se evalúa. Se comparan cuatro métodos, uno de ellos (cero) la referencia, sobre el mismo
sorteo y el mismo resto, de modo que son pareados:

\begin{itemize}
  \item \textbf{Cero.} El modelo \loso{} sin calibrar, medido sobre ese resto. Es la
  referencia pareada: sin ella, la comparación mezclaría el efecto de la calibración con el
  de evaluar sobre menos ensayos.
  \item \textbf{Estadísticos} (\texttt{stats}). Se sustituyen la media y la desviación por canal del
  entrenamiento por las de los $k$ ensayos de calibración. No se modifica ningún peso.
  \item \textbf{Cabeza} (\texttt{cabeza}). Se reentrena solo la última capa de clasificación de la
  arquitectura con los $k$ ensayos, y el resto del modelo permanece congelado. Se emplea
  AdamW con tasa de aprendizaje $10^{-3}$ y decaimiento de pesos $10^{-2}$, entropía cruzada
  sin suavizado y $100$ pasos de lote completo, con las capas de normalización por lotes y
  el \textit{dropout} en modo de evaluación (recalcular sus estadísticos con $10$ ensayos
  sería otra forma de adaptación, no la que se quiere medir). La semilla de la cabeza es la
  semilla de la ejecución más el índice del sorteo.
  \item \textbf{Estadísticos y cabeza} (\texttt{stats+cabeza}). Las dos operaciones: la cabeza se entrena sobre los
  ensayos de calibración ya normalizados con sus propios estadísticos.
\end{itemize}

El número de pasos y la tasa de aprendizaje de la cabeza se fijaron de antemano y no se
ajustaron en ningún conjunto; son una elección del autor, sin una referencia que los respalde, y el
resultado vale para esa configuración. No hay conjunto de validación ni parada temprana en la
calibración. La \emph{ganancia} de un método es su exactitud menos la de \emph{cero} sobre
el mismo sorteo; se promedia entre los cinco sorteos dentro de cada sujeto, y entre
semillas dentro de él.

La calibración se ejecutó sobre los modelos \loso{} de las semillas $2027$ y $2028$ y sobre
los de la variante con alineamiento (semilla $2026$).


\subsection{Indicador de degradación por cambio de dominio}
\label{sec:dsi}

\subsubsection{Problema del indicador clásico}

El indicador de uso común cuantifica la degradación al pasar de la validación
intra-sujeto (\textit{subject-dependent}) a la inter-sujeto
(\textit{subject-independent}), comparación habitual en la literatura de interfaces cerebro-computadora
(BCI) independiente del sujeto \cite{kwon2020}. Esa literatura compara ambos protocolos; la
forma relativa que sigue es la formalización empleada en este trabajo para contrastarla:

\begin{equation}
\mathrm{DSI}_{\text{clásico}} =
  \frac{m_{\text{intra}} - m_{\text{\loso{}}}}{m_{\text{intra}}}
\label{eq:dsi-clasico}
\end{equation}

donde $m$ es una métrica de desempeño. Esta formulación es \textbf{confundida}:
las dos condiciones no difieren solo en si el sujeto de prueba fue visto, sino
también en el tamaño del conjunto de entrenamiento. En \openbmi{}, la condición
intra-sujeto entrena con unos $10^2$ ensayos y la \loso{} con $10^4$. La
diferencia de desempeño resultante mezcla el efecto del cambio de dominio con el
efecto del volumen de datos, y no permite aislar el primero. La
sección~\ref{sec:within-vs-loso} cuantifica esa confusión sobre los datos.

\subsubsection{Indicador propuesto: \texorpdfstring{\dsi{}}{IDS}}

Para evitar la confusión con el indicador clásico, que conserva la sigla \textit{DSI}
establecida en la literatura, el propuesto aquí se denomina \textbf{\dsi{}, índice de
degradación por sujeto}. Se propone comparar, \emph{para un mismo modelo entrenado una sola
vez}, el desempeño sobre dos conjuntos de prueba que difieren únicamente en si el sujeto
que los produjo estuvo o no en el entrenamiento:

\begin{equation}
\mathrm{IDS} = \frac{m_{\text{visto}} - m_{\text{no visto}}}{m_{\text{visto}}}
\label{eq:dsi}
\end{equation}

donde:

\begin{itemize}
  \item $m_{\text{no visto}}$ se evalúa sobre el sujeto excluido de la partición
        \loso{};
  \item $m_{\text{visto}}$ se evalúa sobre una fracción de ensayos
        \textbf{reservada} de los 51 sujetos de entrenamiento, que no participa
        en el ajuste de los pesos, en el cálculo de los estadísticos de
        normalización ni en la parada temprana. Esta última la determina un conjunto
        distinto, el de los dos sujetos de validación, de modo que $m_{\text{visto}}$
        no queda favorecido por el criterio de selección del modelo.
\end{itemize}

La reserva se estratifica por sujeto, clase y sesión, de modo que no se
introduzca cambio temporal (día 1 frente a día 2) en la comparación. Se fijó una
fracción del 10\,\%.

Al compartir modelo, pesos y datos de entrenamiento, ambos conjuntos difieren en la
pertenencia del participante al entrenamiento, y el indicador queda así libre del
confundidor de volumen de datos que afecta a la ecuación~\eqref{eq:dsi-clasico}. No es, sin
embargo, la única diferencia entre ellos. Los ensayos vistos se reservan de las mismas
sesiones que los de entrenamiento, de modo que el modelo se ha ajustado con datos no solo
de esa persona sino también de esa misma sesión de registro, y el indicador refleja ambas
condiciones. Además, el conjunto visto promedia 51 personas y el no visto corresponde a una sola.
Por ello, el indicador de cada \textit{fold} depende casi por completo de la dificultad
del participante excluido, y solo se interpreta de forma agregada. La figura~\ref{fig:particiones} resume el
esquema.

\begin{figure}[H]
\centering
\begin{tikzpicture}[
  font=\footnotesize,
  caja/.style={draw, rounded corners=2pt, minimum height=0.85cm, align=center},
  fl/.style={-{Latex[length=2mm]}, semithick},
  ev/.style={-{Latex[length=2mm]}, semithick, densely dashed, gray!70!black},
  x=1cm, y=1cm
]
\node[caja, fill=black!5, minimum width=9.4cm] (todos) at (0,0)
  {54 sujetos del conjunto de datos};

\node[caja, fill=blue!8,   minimum width=4.6cm] (tr) at (-2.4,-1.7)
  {51 sujetos de entrenamiento};
\node[caja, fill=blue!20,  minimum width=2.0cm] (va) at ( 2.0,-1.7)
  {2 sujetos\\[-3pt] \scriptsize parada temprana};
\node[caja, fill=orange!25, minimum width=2.2cm] (te) at ( 4.6,-1.7)
  {sujeto $s$\\[-3pt] \scriptsize \textbf{no visto}};
\draw[fl] (todos.south -| tr.north) -- (tr.north);
\draw[fl] (todos.south -| va.north) -- (va.north);
\draw[fl] (todos.south -| te.north) -- (te.north);

\node[caja, fill=blue!8,  minimum width=2.6cm] (entren) at (-3.4,-3.4)
  {entrenamiento\\[-3pt] \scriptsize 90\,\% de los ensayos};
\node[caja, fill=green!20, minimum width=1.8cm] (visto) at (-1.0,-3.4)
  {\textbf{visto}\\[-3pt] \scriptsize 10\,\%};
\draw[fl] (tr.south -| entren.north) -- (entren.north);
\draw[fl] (tr.south -| visto.north) -- (visto.north);

\node[caja, fill=black!10, minimum width=2.6cm] (modelo) at (-3.4,-5.1)
  {modelo\\[-3pt] \scriptsize entrenado una vez};
\draw[fl] (entren) -- (modelo);

\draw[ev] (modelo.south) -- (-3.4,-6.2) -- (-1.0,-6.2)
      node[midway, below, font=\scriptsize] {evalúa} -- (visto.south);
\draw[ev] (modelo.south) -- (-3.4,-6.9) -- (4.6,-6.9)
      node[midway, below, font=\scriptsize] {evalúa} -- (te.south);

\node[font=\scriptsize, anchor=west] at (-0.05,-3.4) {\;$\rightarrow m_{\text{visto}}$};
\node[font=\scriptsize, anchor=west] at ( 5.7,-1.7) {\;$\rightarrow m_{\text{no visto}}$};

\node[caja, draw=black, thick, fill=yellow!15, minimum width=6.4cm] (dsi) at (0,-8.3)
  {$\mathrm{IDS} = \dfrac{m_{\text{visto}} - m_{\text{no visto}}}{m_{\text{visto}}}$};
\end{tikzpicture}
\caption{Esquema de particiones para el cálculo del indicador de degradación.
En cada uno de los 54 \textit{folds} \loso{}, un sujeto queda excluido por completo
(\emph{no visto}), dos sujetos completos determinan la parada temprana y se reserva el
10\,\% de los ensayos de los 51 de entrenamiento (\emph{visto}), estratificado por
sujeto, clase y sesión. Ambos conjuntos se evalúan con el mismo modelo y los mismos
pesos, y ninguno de los dos interviene en la selección del modelo.}
\label{fig:particiones}
\end{figure}

\subsection{Criterio operativo de robustez}
\label{sec:criterio-robustez}

El objetivo general de este trabajo plantea la identificación de la arquitectura de mayor
robustez ante el cambio de dominio, lo que exige una definición operativa del término. La
exactitud promedio no la proporciona: las cuatro arquitecturas quedan comprendidas en un
intervalo de $4.6$ puntos en el protocolo \loso{}, y los contrastes pareados de la
sección~\ref{sec:estadistica} no sostienen una diferencia entre las tres primeras.
Calificar como más robusta a la de media más alta equivaldría a atribuir significado a una
diferencia que los datos no respaldan.

Se define aquí la robustez como la estabilidad del desempeño ante desplazamientos del
dominio, y se evalúa sobre tres ejes independientes, medidos con los protocolos descritos
en las secciones precedentes.

El primero es la \textbf{degradación ante participantes no vistos}, cuantificada mediante
el \dsi{} de la sección~\ref{sec:dsi}, que compara un mismo conjunto de pesos sobre
participantes vistos y no vistos. Valores menores corresponden a mayor robustez, con una
salvedad que afecta también al segundo eje: un clasificador en el nivel del azar tendría un
indicador nulo sin ser robusto, de modo que ambos índices se interpretan junto con la
exactitud sobre participantes no vistos.

El segundo es la \textbf{degradación ante un protocolo de adquisición distinto}. Se
introduce para ello un indicador análogo al anterior, el \textbf{\idp{}, índice de
degradación por protocolo}, construido sobre la misma lógica de un único conjunto de pesos
evaluado en dos conjuntos de prueba:

\begin{equation}
\mathrm{IDP} =
  \frac{m_{\text{participante nuevo}} - m_{\text{protocolo nuevo}}}
       {m_{\text{participante nuevo}}}
\label{eq:idp}
\end{equation}

donde $m_{\text{participante nuevo}}$ corresponde a participantes no vistos del conjunto de
origen y $m_{\text{protocolo nuevo}}$ a participantes del conjunto externo. Ambos términos
comparten la condición de no haber sido vistos durante el entrenamiento, de modo que la
diferencia entre ellos no vuelve a contabilizar el cambio de persona que el \dsi{} ya recoge.
Lo que mide es el cambio de protocolo de adquisición en sentido amplio, un factor compuesto
de los elementos de la tabla~\ref{tab:armonizacion}, y no un factor aislado. Su cálculo exige el protocolo de sujetos apartados descrito
en la sección~\ref{sec:protocolos-extra}, pues un modelo entrenado sobre la totalidad de los
participantes no deja ninguno sin ver y no puede producir el primero de los dos términos.
Los dos índices son adimensionales, pero no proceden del mismo régimen: el \dsi{} se calcula
en \loso{} con $51$ sujetos de entrenamiento y $62$ canales, y el \idp{} con un único modelo
de $44$ sujetos y $45$ canales. La comparación de sus magnitudes es, por tanto, aproximada.

El tercero es la \textbf{estabilidad entre participantes}. La exactitud promedio no informa
sobre cómo se reparte el desempeño entre las personas, reparto que resulta determinante
para una interfaz destinada a usuarios arbitrarios. Se cuantifica mediante la desviación
típica de la exactitud por participante, el percentil $5$ de esa distribución y el número
de participantes situados por debajo del umbral funcional del $70$\,\%.

Los tres ejes se reportan de forma conjunta y no se agregan en un índice único. La
agregación requeriría fijar pesos relativos que ninguna consideración del problema
justifica, y encubriría el hecho, relevante en sí mismo, de que las arquitecturas se
ordenan de manera distinta según el eje considerado.

La interpretación del tercer eje admite una ambigüedad que conviene resolver de antemano.
Una arquitectura cuya exactitud media se aproxima al azar comprime necesariamente el rango
de su distribución por participante, de modo que una desviación baja resulta compatible
tanto con un desempeño estable como con una capacidad discriminativa reducida y uniforme.
El percentil $5$ y el recuento bajo umbral separan ambas situaciones, razón por la cual el
eje no se sustenta en la desviación de forma exclusiva. A ello se añade que los tres ejes
se estiman con tres semillas por arquitectura (sección~\ref{sec:semillas-por-arq}), cuya
desviación, calculada con tres valores, es una estimación gruesa, y que el régimen de
optimización no es homogéneo entre esas semillas en \eegnet{}, el Conformer y LaBraM
(sección~\ref{sec:recorte-gradiente}); por ello las diferencias de pequeña magnitud entre
arquitecturas conservan carácter indicativo.

\subsection{Análisis estadístico}
\label{sec:estadistica}

\subsubsection{Métricas}

Se reportan dos. La \textbf{exactitud} es la proporción de ensayos clasificados
correctamente; con dos clases equiprobables, el azar está en 0.50. El
\textbf{coeficiente kappa} \cite{cohen1960} mide lo mismo pero descontando los aciertos
esperables por azar: vale 0 cuando el modelo no aporta nada sobre adivinar y 1
cuando acierta todo. Con dos clases balanceadas en el conjunto de prueba, como es el caso,
kappa es una transformación lineal exacta de la exactitud, $\kappa = 2\,\mathrm{exactitud}-1$,
de modo que ambas métricas ordenan igual y producen los mismos contrastes; que coincidan no
constituye una corroboración. Se reporta kappa porque el \dsi{} calculado sobre él mide la
degradación relativa al exceso sobre el azar, que no depende del nivel base de $0.5$.

\subsubsection{Comparación entre arquitecturas}

Todas las arquitecturas se evalúan sobre \textbf{las mismas 54 particiones}, con
los mismos sujetos en el mismo orden. Esto permite compararlas sujeto a sujeto en
lugar de comparar solo los promedios, que es más informativo: la variabilidad
entre personas es grande y podría impedir detectar una diferencia real entre modelos.

Para decidir si una diferencia observada es atribuible al azar se emplea la
\textbf{prueba de Wilcoxon para muestras pareadas} \cite{wilcoxon1945}, que compara los
resultados sujeto por sujeto sin suponer que sigan una distribución normal. Es bilateral,
sobre los $54$ pares, y reparte las diferencias nulas entre ambos signos
(\texttt{zero\_method="zsplit"} de SciPy). Se acompaña
siempre del número de sujetos en que cada arquitectura resulta mejor, que es una
lectura directa de la consistencia del efecto, y del tamaño de efecto
$r = |Z|/\sqrt{n}$ \cite{fritz2012}, con $Z = \Phi^{-1}(1 - p/2)$ reconstruido a partir del
valor $p$ bilateral y $n = 54$. Las cuatro arquitecturas cuentan con tres
semillas en los dos protocolos. Los contrastes se calculan sobre la media de las tres semillas
($2026$, $2027$ y $2028$) de cada participante, de modo que cada participante aporta un valor por
arquitectura, el número de pares es siempre $54$ y las semillas no se tratan como réplicas
independientes (sección~\ref{sec:estadistica-variantes}); los resultados con la
semilla $2026$ sola y por semilla se informan en la
sección~\ref{sec:variabilidad-semillas}.

Cada protocolo y cada métrica constituyen una familia de seis contrastes, uno por par de
arquitecturas, fijada antes de calcular las cifras con tres semillas. Se reportan los valores $p$
sin corregir y corregidos por Holm \cite{holm1979} dentro de cada familia; con tres semillas la
corrección no cambia ninguna conclusión de significación. La corrección de Holm ordena
los $m$ valores $p$ de menor a mayor, multiplica el $i$-ésimo por $m - i + 1$ e impone que
la secuencia ajustada no decrezca.

\subsubsection{Equivalencia práctica}
\label{sec:equivalencia}

Que una diferencia no alcance significación no demuestra que sea pequeña, ni tampoco lo
demuestra un tamaño de efecto observado reducido. Para sostener que dos arquitecturas rinden
igual se emplea un criterio de equivalencia \cite{schuirmann1987,lakens2017}: se fija un
margen $\delta$, se plantea como hipótesis nula $|\Delta| \geq \delta$ frente a la
alternativa $|\Delta| < \delta$, y se declaran
equivalentes cuando el intervalo de confianza al $90$\,\% de su diferencia pareada queda
íntegramente dentro de $(-\delta,\ +\delta)$, lo que corresponde a dos contrastes unilaterales
al $5$\,\%. El intervalo se obtiene por remuestreo pareado de percentiles
(sección~\ref{sec:bootstrap}) y se aplica a la exactitud por participante, en \loso{} y en
intra-sujeto, calculada como media de las tres semillas de cada participante, de modo que la
variabilidad por inicialización se promedia dentro del participante y no queda incorporada en el
intervalo. Como comprobación se informan también el intervalo jerárquico de sujetos y semillas
(sección~\ref{sec:variabilidad-metodo}), que sí la incorpora, y el contraste TOST con la prueba de
Wilcoxon (corrección de Holm en la familia de seis) y con la $t$ de Student.

Se fija $\delta = 0.03$ de exactitud. Equivale a entre un quinto y un tercio de la
desviación entre participantes, de $0.10$ a $0.15$ según arquitectura y protocolo, y es del
orden de la variabilidad entre semillas de la exactitud media de un participante, de
$0.016$ a $0.025$ en el protocolo intra-sujeto y de $0.019$ a $0.025$ en \loso{}, con las cuatro
arquitecturas y tres semillas. El margen se
fijó después de un primer análisis de los datos y no de forma previa. Por ese motivo, la
sección~\ref{sec:wilcoxon} informa también el menor margen con el que se cumpliría la
equivalencia.

\textbf{Criterios de equivalencia informados.} Con el margen $\delta = 0.03$ de la
sección~\ref{sec:equivalencia} se informan tres criterios, para que la conclusión no
dependa de uno solo: el intervalo al $90$\,\% por remuestreo pareado dentro de
$(-\delta, +\delta)$; la prueba de dos contrastes unilaterales con Wilcoxon, corregida por
Holm en la familia de seis; y la misma con la prueba $t$, sin corregir. Se aplica también al
protocolo intra-sujeto. No se interpreta un resultado no significativo como ausencia de
efecto: se informa como que no se detecta diferencia y, si procede, con el resultado de
equivalencia; y un resultado significativo se acompaña de la magnitud.


\subsubsection{Intervalos de confianza por remuestreo}
\label{sec:bootstrap}

El contraste de Wilcoxon indica si una diferencia es atribuible al azar, pero no cuánta
incertidumbre acompaña a cada cifra reportada. Para acotarla se emplea el
\textbf{remuestreo con reemplazo} (\textit{bootstrap}) \cite{efron1979} sobre los valores por participante
que el entrenamiento ya produjo, procedimiento que no requiere entrenar ningún modelo
adicional.

De los $54$ valores por participante de una arquitectura se extraen, con reemplazo, $54$
valores; se calcula la media de esa extracción; y el procedimiento se repite
$B = 10\,000$ veces. El intervalo de confianza al $95$\,\% son los percentiles $2.5$ y
$97.5$ de esa distribución, sin corrección de sesgo. La semilla del remuestreo se fija en
$2026$, de modo que dos ejecuciones producen cifras idénticas; como cada cantidad consume
el generador en un orden distinto, la referencia para reproducir una cifra concreta es la
función del verificador que la calcula. El remuestreo trata los \textit{folds} \loso{} como
independientes, aunque comparten $49$ de sus $51$ sujetos de entrenamiento.

La aplicación del procedimiento al \dsi{} requiere una precisión adicional. El remuestreo
se realiza sobre los \textit{folds} y el indicador se recalcula en cada extracción con la
ecuación~\eqref{eq:dsi} aplicada a las medias, que es la forma en que se reporta; promediar
el indicador de cada \textit{fold} estimaría una cantidad distinta. En las diferencias
entre arquitecturas, por su parte, la extracción de participantes es común a las dos
arquitecturas del par, pues remuestrear cada una por separado destruiría la correlación
entre ellas e inflaría el intervalo de forma artificial.

Con tres semillas se añade un remuestreo jerárquico. En cada una de las $B = 10\,000$ repeticiones
se extraen con reemplazo $54$ participantes, dentro de cada uno se extraen con reemplazo sus tres
semillas, con índices independientes entre arquitecturas, y se promedian las semillas antes de
calcular la diferencia o el indicador. Remuestrear en los dos niveles cuenta dos veces el ruido de
semilla, que ya forma parte de la dispersión entre participantes, de modo que estos intervalos son
algo conservadores; con solo tres semillas son además una aproximación tosca de la variación por
inicialización. Los intervalos de remuestreo solo de participantes difieren de los jerárquicos en
menos de $0.003$ en las diferencias entre arquitecturas.

Los intervalos no incorporan corrección por comparaciones múltiples. Se reportan como
descripción de la incertidumbre de cada cifra; las afirmaciones sobre diferencias entre
arquitecturas se sustentan en el contraste pareado, con la comprobación de Holm descrita
arriba.

\subsubsection{Criterio para declarar degradación}

Se declara degradación por cambio de dominio únicamente si se cumplen dos
condiciones: que el \dsi{} sea positivo y que la diferencia entre el conjunto
visto y el no visto no sea atribuible al azar, según un contraste de Wilcoxon bilateral
entre la exactitud en vistos y en no vistos sobre los $54$ \textit{folds}, con
$\alpha = 0.05$. Si solo se cumple la primera, se
reporta el valor indicando expresamente que no es concluyente.

Estas dos condiciones se fijaron antes de ejecutar las arquitecturas restantes. La
corrección de Holm sobre los cuatro contrastes, uno por arquitectura, se incorporó
\textbf{después} de conocer los resultados de las cuatro, y determina la clasificación de una
de ellas: sin corrección, EEGMamba ($p = 0.038$) también se declararía degradada. Por eso
la sección~\ref{sec:dsi-resultados} informa ambos valores.

\subsubsection{Contrastes de las variantes y familias de comparaciones múltiples}
\label{sec:estadistica-variantes}

\textbf{Unidad y contraste.} La unidad de análisis es el sujeto ($n = 54$). En el protocolo
intra-sujeto, la exactitud de un sujeto en una semilla es la media de sus $20$ particiones;
en \loso{} es la de su \textit{fold}. Cuando hay varias semillas, se promedian dentro del sujeto o
se modelan como factor anidado; nunca se tratan como réplicas independientes. El contraste
de una variante es la diferencia por sujeto entre la variante y su referencia de la misma
semilla. Para cada contraste se informa la media y la mediana de la diferencia, el
intervalo al $95$\,\% por remuestreo de sujetos (percentiles, $B = 10\,000$, con un
generador \texttt{numpy.random.default\_rng(2026)} nuevo para cada contraste), la prueba de
Wilcoxon bilateral con reparto de las diferencias nulas, el tamaño de efecto $r$ y el número
de sujetos que mejoran y empeoran, como en los contrastes entre arquitecturas descritos arriba. Dado que kappa
es una transformación lineal de la exactitud, se contrasta la exactitud.

\textbf{Familias de comparaciones múltiples.} La corrección de Holm se aplica dentro de
cada familia, definida antes de calcular los contrastes de las variantes y de las semillas adicionales (resumen en la tabla~\ref{tab:familias}):

\begin{itemize}
  \item entre arquitecturas, seis contrastes por protocolo, métrica y variante de semilla
  (semilla $2026$, media de tres semillas, media de $2027$ y $2028$);
  \item ablaciones intra-sujeto, una familia por arquitectura: ocho contrastes para
  \eegnet{} y para el Conformer, y tres para LaBraM y para EEGMamba;
  \item \loso{} con alineamiento: cuatro contrastes, uno por arquitectura; \loso{} con
  reflexión de canales: cuatro, uno por arquitectura;
  \item calibración con pocos ensayos: veinticuatro contrastes de ganancia frente a cero
  (cuatro arquitecturas, dos valores de $k$, tres métodos), para los modelos base y,
  aparte, otros veinticuatro para los de la variante con alineamiento; el contraste de
  equivalencia de cada una forma una familia propia de veinticuatro;
  \item comparación de las políticas de normalización del destino: ocho contrastes;
  \item alineamiento en la transferencia: doce contrastes (tres evaluaciones por cuatro arquitecturas).
\end{itemize}

Como lectura secundaria declarada se aplica Holm sobre los $22$ contrastes de las
ablaciones a la vez y sobre los $8$ de \loso{} con alineamiento y con reflexión a la vez. Como
análisis de sensibilidad, se repiten las familias de alineamiento, reflexión y las
ablaciones de LaBraM y EEGMamba con la referencia sustituida por la media por sujeto de las
semillas $2027$ y $2028$ de la base, entrenadas con recorte; para \eegnet{} y el Conformer en
intra-sujeto no existe esa referencia, porque sus semillas $2027$ y $2028$ se entrenaron sin
recorte (tabla~\ref{tab:regimen}). Todo análisis fuera de estas familias es exploratorio y no se
corrige.


\subsubsection{Variabilidad por semilla de las diferencias}
\label{sec:variabilidad-metodo}

\textbf{Variabilidad por semilla de las diferencias.} Para cada contraste entre
arquitecturas se calcula la diferencia media por semilla (emparejando semillas por índice) y
la desviación típica de esos tres valores, con la salvedad de que tres valores dan dos
grados de libertad. Además, el intervalo de la diferencia de las medias de tres semillas
se obtiene con el remuestreo jerárquico descrito en la sección~\ref{sec:bootstrap}, con un
generador \texttt{default\_rng(2026)} nuevo por contraste. El mismo procedimiento, con los \textit{folds} \loso{}
como unidad en el \dsi{} y con los sujetos apartados y los de WBCIC-SHU como unidades del
\idp{}, da los intervalos jerárquicos de ambos índices.


\subsection{Control de fuga de datos}

El particionado \loso{}, el intra-sujeto y el conjunto de vistos están cubiertos por
\textbf{41 pruebas automatizadas}: 22
sobre el particionado \loso{} y la separación de sujetos apartados, 12 sobre el
intra-sujeto y 7 sobre el conjunto reservado. No las tienen, en cambio, las funciones que
construyen el protocolo global, el de sujetos apartados y el $20$\,\% reservado de
WBCIC-SHU; para esos tres se comprobó que los recuentos de segmentos coinciden con los
esperados. Las pruebas verifican, entre otras condiciones, que ningún ensayo reservado
participe en el entrenamiento, que provenga de sujetos de entrenamiento, que
esté estratificado por clase y sesión, y que los estadísticos de normalización
se ajusten exclusivamente sobre el conjunto de entrenamiento reducido.

\subsection{Cota del efecto del cambio de régimen}
\label{sec:cota-confusor}

Para acotar cuánto del cambio entre la semilla $2026$ y las semillas $2027$ y $2028$ puede
deberse al régimen y no al azar, se calcula para cada sujeto $d_{\mathrm{cruce}} =
\mathrm{exac}(2026) - \tfrac12\bigl(\mathrm{exac}(2027) + \mathrm{exac}(2028)\bigr)$, con la
exactitud media del sujeto en cada semilla, y $d_{\mathrm{intra}} = \mathrm{exac}(2027) -
\mathrm{exac}(2028)$. Si el régimen no importara y el ruido de semilla tuviera la misma
varianza en las tres, ambas diferencias tendrían media cero, y la varianza de
$d_{\mathrm{cruce}}$ sería $0.75$ veces la de $d_{\mathrm{intra}}$. EEGMamba, que no cambia de
recorte, actúa como control; en el protocolo intra-sujeto, \eegnet{} y el Conformer, que
nunca tuvieron recorte, son un control sin cambio de recorte ni de hardware entre
semillas.

La cota del efecto sobre la exactitud media por sujeto es el mayor valor absoluto de los
extremos del intervalo al $95$\,\% (percentiles de un remuestreo de sujetos, $B = 10\,000$) de
la media de $d_{\mathrm{cruce}}$, o de su diferencia con la de EEGMamba si esa es mayor. En
una media de tres semillas, de las cuales una se entrenó sin recorte, la contaminación
máxima es la tercera parte de la cota. El análisis supone aditividad (el cambio de régimen
suma un desplazamiento constante a la exactitud de cada sujeto) y ruido de semilla
independiente; el cambio de régimen es compuesto (recorte, hardware y código) y no puede
separarse; y el control de EEGMamba supone que su sensibilidad al hardware y al código es la
de las demás arquitecturas, lo cual no es verificable. Es un análisis exploratorio y sus
pruebas no se corrigen por comparaciones múltiples.

El hardware de ambos regímenes difiere (tabla~\ref{tab:hardware}); no se interpreta ninguna
diferencia como efecto del hardware.


% =====================================================================
%  RESULTADOS
% =====================================================================
